\chapter{Conclusions and Future Work}
\label{chap:conclusions-future-work}

This chapter concludes the research and development documented in this thesis. The primary problem addressed by this work is the growing social isolation and lack of mutual aid within high-density urban apartment communities, exacerbated by the absence of trusted, hyper-local digital communication channels. The goal was to engineer a structured, secure web platform that bridges this social gap. The proposed solution, NEST, is a hyper-local web application that provides modules for community feeds, peer-to-peer tool borrowing, parking spot sharing, and event organization. 

The main advantages of the NEST platform include its strict building-level verification, which ensures a high-trust environment, and its highly responsive, reactive frontend architecture that provides an excellent user experience. However, the current implementation has disadvantages, such as the reliance on an SQLite database that limits horizontal scaling, and the absence of real-time WebSocket capabilities, meaning users must manually refresh to see immediate updates.

\section{Summary of Contributions}
By analyzing the system against the core objectives defined in Chapter 1, all major milestones have been achieved:
\begin{itemize}
    \item \textbf{Milestone 1: Security \& Verification (O1):} Successfully implemented a multi-stage onboarding flow. Guests are verified by local administrators using secure blocks, restricting data access to verified residents.
    \item \textbf{Milestone 2: Cohesive Design System (O2):} Created a design token architecture with over 600 CSS variables. This supports light and dark themes and maintains responsive layouts across mobile, tablet, and desktop viewports.
    \item \textbf{Milestone 3: Collaborative Consumption (O3):} Developed the Shared Shed module, providing a digital tool catalog with conflict prevention logic and state-machine transitions to support resource sharing.
    \item \textbf{Milestone 4: Parking Optimization (O4):} Built the Parking Sharing module. This includes an automated conflict-resolution engine that denies overlapping requests once an application is approved, ensuring efficient spot booking.
    \item \textbf{Milestone 5: High-Performance Reactive Architecture (O5):} Integrated the NgRx Facade state model with Angular Signals. This represents an innovative adaptation of the traditional NgRx Facade pattern---originally built for RxJS Observables---newly retrofitted to support Angular's modern Signal APIs. This isolates store operations, handles async data pre-fetching, and enables fine-grained DOM rendering to minimize browser performance overhead.
\end{itemize}

\section{Lessons Learned and Engineering Limitations}

Developing NEST provided several critical software engineering insights. Transitioning from traditional Angular change detection to reactive Signal APIs simplified component state management. Monotonically tracking inputs and computed values reduced boilerplate lifecycle code while improving rendering precision. Furthermore, implementing a Facade layer between the visual components and the NgRx store vastly improved maintainability. Components remain decoupled from raw actions and selectors, allowing development teams to refactor state schemas without altering template markup. On the styling front, centralizing the design system inside a single, structured color token sheet simplified layout development, proving that robust CSS architectures prevent styling duplication and keep components theme-agnostic.

Despite these successes, several limitations were identified during the development and testing phases. All updates, such as feed posts and shed requests, currently rely on standard HTTP requests and store refreshes. The lack of WebSocket integration prevents true real-time notifications, requiring manual page refreshes or periodic polling. On the backend, the SQLite database is excellent for local staging and single-host deployments, but it lacks the concurrent transaction scaling needed for multi-block, high-throughput enterprise environments. Additionally, profile and post image attachments are saved directly to local directory systems on the API server, which lacks the durability and horizontal scaling provided by cloud object stores.

\section{Future Work}

To address these limitations and expand the platform's capabilities, several future enhancements are planned. The most immediate priority is the integration of WebSocket real-time notifications using Socket.io on the Express backend and RxJS WebSocket subjects on the Angular client. This will enable instant, real-time chats, active borrowing requests, and parking alerts without manual refreshes.

To prepare the application for enterprise production, the persistence layer will be migrated to PostgreSQL to support concurrent queries, and static uploads will be transitioned to secure Amazon S3 or Google Cloud Storage buckets. Finally, configuring Progressive Web Application (PWA) capabilities through service workers will support offline caching and native push notifications on mobile devices, while integrating Angular's localization libraries will provide support for multiple languages, broadening the application's accessibility.

In conclusion, NEST demonstrates how structured, hyper-local platforms can actively support mutual aid and community connection. The architectural patterns, reactive flows, and security designs implemented in this project provide a stable and highly maintainable foundation for future local community engineering.
